% !PRINTED 148
\weberpage{148}
The degree $g_2$ is obtained exactly like that of $Q_0$ and $Q_{0,1}$. But if
under $Q_{0,1,2}$ a fourth figure 3 were also to remain unchanged, then $g_2$
would be a divisor of the degree of the symmetric group of $n-4$ figures, hence
a divisor of $\Pi(n-4)$. Thus

\[ \frac{n\,\Pi(n-4)}{\Pi(n-3)} = \frac{n}{n-3} \]

would have to be an integer, hence $n \geqq 2n - 6$ or $n \leqq 6$.

From $\varepsilon$) we conclude that there is in $Q$ a permutation $\varkappa$
by which $0, 1, 2, 3$ passes into $0, 1, 2, 4$, where 4 is a figure different
from 3. By $\beta$), then, $Q$ also contains the permutation:

\[ \varkappa^{-1}(0,1)(2,3)\varkappa = (0,1)(2,4), \]

and consequently also:

\[ (0,1)(2,3)(0,1)(2,4) = (2,3,4), \]

that is, a three-term cycle, which contradicts the hypothesis. Hereby
Theorem II. is completely proved.

That the cases $n = 4$ and $n = 6$ really form exceptions follows from Bd. I,
\S.\,160 and \S.\,182, where we have seen \gesperrt{that the symmetric permutation group
of four figures possesses a transitive divisor of index 3 and that of six
figures a transitive divisor of index 6.}

\begin{center}
\rule{3cm}{0.4pt}
\end{center}
